\section{Описание практической части}
\label{sec:Chapter5} \index{Chapter5}

%Если в рамках работы писался какой-то код, здесь должно быть его
%описание: выбранный язык и библиотеки и мотивы выбора, архитектура,
%схема функционирования, теоретическая сложность алгоритма, характеристики
%функционирования (скорость/память).

Для написания кода использовался язык \texttt{Python} и опирающийся на его синтаксис 
\texttt{SageMath}. В качестве криптографических протоколов, основываясь на статье~\cite{RLWE_Channel}, были выбраны $\LAC$ и $\NewHope$.

Зафиксированы следующие наборы параметров:
\begin{table}[h]
	\centering
	\begin{tabular}{|c|c|c|c|}
		\hline Протокол & n & q & k \\ \hline
		$\LAC$ & $1024$ & $251$ & $1$ \\ \hline
		$\NewHope$ & $1024$ & $12289$ & $8$ \\ \hline
	\end{tabular}
\end{table}

\vspace{-0.5cm}
\subsection{Вычисление шума}

Вероятности исходов $\mathcal{B}(k, \frac12)$ посчитаны по формулам через модуль \texttt{scipy.stats}. После этого распределение $\chi_k$ вычислено за время $O(k^2)$ по определению: перебором возможных значений $(\mu, \nu)$, вычислением значения $\mu-\nu$ и прибавлением к нему вероятности текущего исхода.

В распределении $\mathcal{B}(k, p)$ присутствуют значения $0 \leq x \leq k$, отсюда в $\chi_k$ они будут $-k \leq x \leq k$ --- всего $O(k)$ штук. Таким образом, $\xi_k$ также вычисляется по определению за $O(k^2)$ тем же способом.

После этого нам нужно свернуть $\xi_k$ само с собой $2n$ раз, а после этого ещё $1$ раз с $\chi_k$ для получения распределение шума $\psi$. Все операции происходят в $\Z_q$, поэтому на каждом шаге распределения принимают не больше $q$ разных значений $0 \leq x < q$.

Сворачивать с собой $\xi_k$ будем, используя алгоритм бинарного возведения в степень и тратя на это $O(\log(n))$ шагов. Во время каждой его итерации, а также в конце при свёртке с $\chi_k$ нам нужно находить свёртку распределений со значениями из $\{0, \ldots, q-1\}$.

Обозначая их $P$ и $Q$, значения $P * Q$ ищутся в $\Z_q$ по формуле:
\begin{equation*}
	(P*Q)[y] = \sum_{x=0}^y{P[x]Q[y-x]}+\sum_{x=y+1}^{q-1}{P[x]Q[q+y-x]}
\end{equation*}

Если взять проассоциированные с распределениями многочлены, свёртка выражается через коэффициенты их произведения:
\begin{equation*}
	(P*Q)[y] = (\hat P\hat Q)[y]+(\hat P\hat Q)[q+y]
\end{equation*}

Но произведение многочленов длины $q$ можно с помощью $\mathsf{NTT}$ искать за время $O(q\log(q))$. В моём коде для этих целей я использую функцию \texttt{fftconvolve} из модуля \texttt{scipy.signal}.

В итоге, распределение $\psi$ может быть посчитано с использованием $O\Bigl(q\log(q)\log(n)+k^2\Bigr)$ времени и $O\Bigl(q\log(n)\Bigr)$ памяти, чего вполне хватает для большинства используемых сейчас криптосистем.

\subsection{Вероятности ошибок}

Дальше, для быстрого вычисления $\Demap$ предподсчитаем его значения. Будем перебирать символ $x \in \Z_q$ и смотреть, к образу какого из $s \in \Sigma_Q$ он ближе --- асимптотика $O(q\?Q)$.

Теперь $\Map$ и $\Demap$ считаются за $O(1)$, и можно по определению вычислить значения $ch$ --- также за $O(q\?Q)$.

Зная $\psi$, по ним за $O(q)$ восстанавливается распределение $\sigma$. Наконец, для оценки $\DFR$ у нас снова возникает задача посчитать $P^n$, но уже в $\R[x]$.

Посмотрим на устройство алгоритма бинарного возведения в степень. Так как $\deg(P) = Q-1$ и, следовательно, $\deg(P^m) \leq (Q-1)\?m$, можно дополнять $P^m$ нулями до длины $Qm$.

В процессе вычисления мы $\leq 2$ раза перемножим многочлены длины $\frac12 Qn$, $\leq 2$ раза --- длины $\frac14 Qn$ и т.д. С $\mathsf{NTT}$ суммарно получается:
\begin{equation*}
	Time \leq \sum_{i=1}^B{\frac1{2^i}Qn\cdot\log\left(\frac1{2^i}Qn\right)}
\end{equation*}

Огрубляя и оценивая сумму геометрической прогрессии, выводим:
\begin{equation*}
	Time \leq Qn\log(Qn)
\end{equation*}
Имеем асимптотику $O\Bigl(Q(n\log(Qn)+q)\Bigr)$ времени и $O\Bigl(q+Qn\Bigr)$ памяти.

\subsection{Подсчёт t и R}

\begin{state}\label{st:search}
	Значение $\DFR$ при фиксированном $n$ убывает при возрастании $t$, из-за этого первый момент выполнения $\DFR < 10^{-6}$ можно искать с помощью алгоритма бинарного поиска.
\end{state}

Вычисление $\DFR$ происходит суммированием за $O(Qn)$, значение $t$ ограничено тем же числом --- асимптотика $O(Qn\log(Qn))$.

Дальше по данным $t$ и $n$ нужно вычислить $M$. В случае Хэмминга я применяю для этого методы \texttt{gilbert\_lower\_bound} и \texttt{hamming\_upper\_bound} из библиотеки \texttt{sage.codes.bounds}.

В метрике Ли я явно строю $P$, как элемент \texttt{PolynomialRing(ZZ)} из \texttt{SageMath}, и после встроенным методом считаю $P^n$, откуда узнаю объёмы шаров и границы на $M$ соответственно.

Данный шаг снова работает за асимптотику $O(Qn\log(Qn))$. Наконец, чтобы узнать $R$, я рассматриваю $M$, как элемент \texttt{RealField(100)} из \texttt{SageMath}, и встроенной функцией нахожу его логарифм.
